\documentclass[11pt]{article}
\usepackage[letterpaper,left=1in,right=2.3in,top=1in,bottom=1in,
            marginparwidth=1.8in,marginparsep=0.2in]{geometry}
\usepackage{amsmath,amssymb,amsthm}
\usepackage{stmaryrd}
\usepackage{mathpartir}
\usepackage{xcolor}
\usepackage{booktabs}
\usepackage{array}

\newtheorem{definition}{Definition}
\newtheorem{conjecture}{Conjecture}
\newtheorem{lemma}{Lemma}
\newtheorem{theorem}{Theorem}
\theoremstyle{remark}
\newtheorem*{remark}{Remark}

% review notes for the author
\newcommand{\todo}[1]{\marginpar{\raggedright\footnotesize\textcolor{red!70!black}{\textbf{TODO:} #1}}}

% languages
\newcommand{\lpr}{\lambda_{\mathrm{PR}}}
\newcommand{\Nat}{\mathsf{Nat}}
\newcommand{\Unit}{\mathsf{Unit}}
\newcommand{\tRef}{\mathsf{Ref}\,}
\newcommand{\Type}[1]{\mathsf{Type}_{#1}}
\newcommand{\arr}[1]{\xrightarrow{#1}}
\newcommand{\new}{\mathsf{new}\,}
\newcommand{\deref}{!\,}
\newcommand{\assign}{\mathrel{:=}}
\newcommand{\unit}{\langle\rangle}
\newcommand{\lam}[2]{\lambda #1{:}#2.\,}
\newcommand{\klet}{\mathsf{let}\ }
\newcommand{\kin}{\ \mathsf{in}\ }
\newcommand{\Prom}[2]{\lceil #2 \rceil^{#1}}
\newcommand{\prom}[1]{\lceil #1 \rceil}
\newcommand{\frc}[1]{\lfloor #1 \rfloor}
\newcommand{\snat}{\mathsf{nat}}
\newcommand{\slam}[1]{\lambda #1.\,}
\newcommand{\tr}[1]{\llbracket #1 \rrbracket}
\newcommand{\ext}{\sqsupseteq}
\newcommand{\Heap}{\mathit{Heap}}
\newcommand{\Comp}{\mathit{Comp}}
\newcommand{\Val}{\mathit{Val}}
\newcommand{\Loc}{\mathit{Loc}}

\title{Sharing is not heap-adequately expressible\\
       in predicative type theory}
\author{}
\date{Draft}

\begin{document}
\maketitle

\begin{abstract}
  Call-by-need STLC is total, and pure promises are harmless in type theory:
  a promise can be read as a function from $\Unit$. Yet we have been unable
  to give call-by-need STLC a shallow embedding in a predicative type theory
  in which the heap is a semantic object. We argue that this is not an
  accident of one construction. A semantic heap of thunks is an instance of
  the type-world circularity, and a single program, $P = \slam{g} g\,(\slam{n}
  g\,(\slam{m} m))$, defeats every well-founded stratification of it. We
  make the claim precise as a failure of \emph{heap-adequate
  expressibility}: compositional translations, required to preserve types
  and to preserve the heap. Unlike Felleisen's macro-expressibility, the
  translation may reinterpret every construct, so store-passing is allowed;
  like it, each construct must map to a fixed context. Target cells must correspond one
  to one with live source cells, with related contents. Heap adequacy is how
  the definition sees \emph{sharing}: a translation that recomputes instead
  of sharing needs either a missing cell or two cells where the source has
  one. Under that definition, $\lpr$'s universe levels are exactly the price
  of embedding stateful programs in type theory, and the level-free promise
  calculus is more expressive than predicative type theory.
\end{abstract}

\section{Introduction}

The goal that started this work was a shallow embedding of call-by-need STLC
in Agda: types as Agda types, terms as Agda functions, a heap holding
semantic values, and totality by structural recursion alone. The heap is not
observable by programs. We want it because it determines \emph{cost}: what is
allocated, what is forced, what is shared, and what becomes garbage.

A stratified model, in the style of $\lpr$~\cite{koronkevich2025}, gets
partway. It is shallow and total, but it covers only a fragment of STLC. The
program
\[
  P \;\triangleq\; \slam{g} g\,(\slam{n} g\,(\slam{m} m))
  \;:\; ((\snat \to \snat) \to \snat) \to \snat
\]
is outside the fragment, and no choice of levels admits it. A logical
relation over syntactic stores does cover all of STLC, and proves
termination with no levels at all, but it puts syntax in the heap.

This note asks what that gap means. Our answers:
\begin{itemize}
  \item The circularity a shallow heap meets is the type-world circularity,
    with the heap's type playing the semantic world (\S\ref{sec:circ}).
  \item $P$ defeats every well-founded measure on that circularity, not just
    universe levels (\S\ref{sec:obstruction}).
  \item The gap is invisible to observational notions of expressiveness. A
    heap-sensitive notion, which we call heap-adequate expressibility,
    captures it. Unlike Felleisen's macro-expressibility, it allows
    store-passing translations, which both sides of our comparison need.
    The feature it isolates is sharing, which is what makes call-by-need
    cheaper than call-by-name (\S\ref{sec:cost}, \S\ref{sec:def}).
  \item Under that notion we conjecture that $\lpr$ is expressible in
    predicative type theory and call-by-need STLC is not
    (\S\ref{sec:conj}). The dividing line is the compositionality
    condition: a translation may use the translations of subterms, but not
    their syntax.
\end{itemize}
Companion artefacts:
\texttt{counterexample.tex} has the derivations and the logical relation in
detail; \texttt{Stratified.agda} is the stratified shallow model; and
\texttt{Promises.agda} mechanizes the syntactic model.

\section{Background}

\subsection{$\lpr$}

$\lpr$ is a call-by-value calculus with higher-order references and a
hierarchy of universes~\cite{koronkevich2025}. A reference costs a level:
$\tRef \tau :: \Type{i+1}$ when $\tau :: \Type{i}$. A function type carries a
level, $\tau_1 \arr{k} \tau_2$, and a $\lambda$ at level $k$ requires $k$
to bound the levels of $\tau_1$, $\tau_2$ and the context. The bound is
non-strict, because functions are immutable. Each universe is a Kripke
world: a world at level $k+1$ maps locations to semantic types at levels
$\le k$, which breaks the type-world circularity.

\paragraph{The $\lambda$ rule.}
The paper's main typing rule (its Fig.~3) bounds the whole context,
$k \geq \max(\Gamma, \tau_1, \tau_2)$. Its \S3.1 gives a variant that bounds
only the free variables, $k \geq \max(\mathit{FV}(e), \tau_1, \tau_2)$, which
also restores weakening and substitution. We use the free-variable rule
throughout. It is the more realistic choice for an implementation, and it
matters in practice. With the whole-context rule, a closed function written
inside the scope of a reference, such as $\lambda x. 7$ under
$r : \tRef (\Nat \arr{0} \Nat)$, is forced to level~$1$ and cannot be stored
in $r$. In the mechanization, the body of a $\lambda$ is typed in a
sub-context $\Delta \subseteq \Gamma$ given by a thinning, and $k$ bounds
$\Delta$. Choosing $\Delta = \mathit{FV}(e)$ gives the free-variable rule.
\todo{Check the statement of the free-variable $\lambda$ rule against \S3.1 of the paper.}

\subsection{Call-by-need as promises}

A promise is a memoizing thunk. Encoded in $\lpr$:
\[
\begin{array}{@{}l@{\;\;\triangleq\;\;}l@{}}
  \Prom{k}{\tau} & \tRef (\Unit \arr{k} \tau) \\
  \prom{e} & \new (\lam{\_}{\Unit} e) \\
  \frc{x} & \klet v = (\deref x)\,\unit \kin x \assign (\lam{\_}{\Unit} v);\ v
\end{array}
\]
so that $\Prom{k}{\tau} :: \Type{k+1}$. The translation of call-by-need STLC
makes every variable and every argument a promise:
\[
\begin{array}{@{}l@{\;\;=\;\;}l@{\qquad\qquad}l@{\;\;=\;\;}l@{}}
  \tr{\snat} & \Nat & \tr{x} & \frc{x} \\
  \tr{A \to B} & \Prom{j}{\tr{A}} \arr{k} \tr{B} &
  \tr{\slam{x{:}A} e} & \lam{x}{\Prom{j}{\tr{A}}} \tr{e} \\
  \multicolumn{2}{l}{} & \tr{e_1\,e_2} & \tr{e_1}\,\prom{\tr{e_2}}
\end{array}
\]

\subsection{The witness}

The image of $P$ is
$\lam{g}{T} \frc{g}\,\prom{\lam{n}{N} \frc{g}\,\prom{\lam{m}{N} \frc{m}}}$,
with $N = \Prom{0}{\Nat}$, $F = N \arr{c} \Nat$, $R = \Prom{b}{F}$,
$\tau_g = R \arr{d} \Nat$ and $T = \Prom{a}{\tau_g}$. Kinding gives
$c \geq 1$, $b \geq c$, $d \geq b+1$ and $a \geq d$. The closure
$\lam{n}{N} \cdots$ captures $g$, so $c \geq a+1$. Together:
\[
  c \;\geq\; a+1 \;\geq\; d+1 \;\geq\; b+2 \;\geq\; c+2.
\]
Both $+1$s are the $\mathsf{Ref}$ inside a promise: $g$'s argument is a
cell, and the closure captures $g$'s cell. Read as a plain CBV term, $P$
types with every level at $0$: the same cycle, with every edge non-strict.
Passing values instead of promises where possible does not help: when $P$ is
applied to a computation, $g$'s cell is unavoidable, and the chain becomes
$c \geq a+1 \geq d+1 \geq c+1$.

Running $P$ on $\slam{f} f\,0$ allocates a cell $g$, then a cell $f_1$
holding $\slam{n} g\,(\slam{m} m)$, which captures $g$ and is passed to $g$'s
function. The heap at run time is acyclic: $f_1$ points to $g$, which points
nowhere. But $g$ must sit above $f_1$, because $g$'s function is applied to
$f_1$, and $f_1$ must sit above $g$, because it captures $g$.

\section{The shallow heap is the type-world circularity}
\label{sec:circ}

A shallow model of call-by-need with a heap has the shape
\[
\begin{array}{@{}l@{\;\;=\;\;}l@{}}
  \Comp\,\tau\,W & \Heap\,W \to
    \Sigma\,W' \ext W.\ \Heap\,W' \times \Val\,\tau\,W' \\
  \Heap\,W & \prod_{(\ell : \sigma) \in W}
    \Val\,\sigma\,W \uplus \Comp\,\sigma\,W \\
  \Val\,(A \to B)\,W & \forall W' \ext W.\
    \Loc\,W'\,A \to \Comp\,B\,W'
\end{array}
\]
$\Heap$ occurs negatively in its own definition. Even when worlds $W$ are
syntactic lists of types, $\Heap\,W$ is the interpretation of $W$, cell by
cell. That is Ahmed's equation~\cite{ahmed2004}, with the heap's type
playing the semantic world:
\[
\begin{array}{@{}l@{\qquad\sim\qquad}l@{}}
  \mathit{Type} = \mathit{World} \to \mathcal{P}(\Val)
    & \Comp\ \text{reads}\ \Heap \\
  \mathit{World} = \Loc \rightharpoonup \mathit{Type}
    & \Heap\,W = \textstyle\prod_\ell \llbracket W(\ell) \rrbracket
\end{array}
\]
\todo{Cite Ahmed's thesis (or Ahmed, Appel and Virga) for the type-world circularity; confirm the right reference.}
A predicative type theory accepts such a definition only with a well-founded
measure that decreases from a cell's domain to the heaps that domain reads.
$\lpr$'s levels are such a measure, and so are Agda's universe levels: a
stratified model puts only $\Val_k$ and $\Comp_k$ in $\Heap_{k+1}$.

The four models of higher-order state we know of avoid the circularity in
four ways:
\begin{itemize}
  \item $\lpr$ stratifies worlds by universe level.
  \item Extended $L^3$ consults its store typing only at a smaller
    step index~\cite{ahmed2007}.
  \item Full ground references restrict contents to first-order data, so
    the heap's interpretation never reads a heap~\cite{kammar2017}.
  \item The syntactic model of promises never forms the heap's
    interpretation at all: cells hold syntax, and the meaning of a cell is
    carried by the type of the pointer to it.
\end{itemize}

\section{The obstruction}
\label{sec:obstruction}

\paragraph{The witness defeats every measure.}
Run $P$ on $\slam{f} f\,0$, and write $D(x)$ for the type interpreting the
contents of cell $x$.
\begin{itemize}
  \item $g$ holds a function that forces its argument $f_1$, so $D(g)$ takes
    heaps containing $f_1$: $\mu(g) > \mu(f_1)$.
  \item $f_1$ holds a thunk that forces $g$, so $D(f_1)$ takes heaps
    containing $g$: $\mu(f_1) > \mu(g)$.
\end{itemize}
Both inequalities come from the definitions of the domains, not from any
particular run.

\paragraph{Capture versus use.}
The history of a run is stratified: every pointer out of a cell points into
that cell's past, if cells allocated while forcing a cell are ordered just
before it. But that orders what a cell's contents \emph{capture}. Defining
$D$ needs what they \emph{do}, and functions are applied in the future. $g$'s
function captures nothing, yet it is applied to $f_1$, which is created after
$g$ and reads $g$. Interpreting $g$'s contents in a world of cells up to $g$
means $g$'s function can only be applied to cells up to $g$, which is
$\lpr$'s restriction again.

\paragraph{Fixes that fail.}
\begin{itemize}
  \item \emph{Measure by type.} $f_1 : F$ reads $g : \tau_g$, a larger type.
  \item \emph{Measure by allocation time or dependency position.} These grow
    with the history, and the future is unbounded.
  \item \emph{Levels chosen at run time.} $g$'s value, and so its level, is
    fixed before $f_1$ exists, and $f_1$ must sit above $g$'s cell.
  \item \emph{Separate structures for forced and unforced cells.} The cycle
    passes through both: when $g$'s function is applied, $g$ is forced and
    $f_1$ is not.
  \item \emph{Level-polymorphic stored functions.} This is impredicative: in
    Agda such a value lives in $\mathsf{Set}\omega$, and no cell can hold it.
    An impredicative universe does not help, because the equation is not
    positive, and impredicativity is inconsistent with non-strictly-positive
    inductive types~\cite{coquand1990}.
    \todo{Confirm the Coquand--Paulin-Mohring result and its exact statement (positive but not strictly positive types with impredicative Prop).}
  \item \emph{One copy of $g$ per use.} That is call-by-name; \S\ref{sec:cost}
    shows the cost can be exponential.
\end{itemize}

\paragraph{Execution is well-founded; definition is not.}
Forcing a cell runs only on its dependency prefix and the cells it
allocates. Even the inner call of $g$'s function, inside the forcing of
$f_1$, happens in a world without $f_1$. But $D(g)$ must also cover the outer
call, in a world containing $f_1$. The syntactic model exploits the
well-founded execution and never defines $D$.

\section{Why the heap matters: cost}
\label{sec:cost}

The heap is unobservable, so any notion of expressiveness that looks only at
results sees nothing. But the observationally correct translations of
promises are wrong for cost:
\begin{itemize}
  \item \emph{Call-by-name}, $\prom{e} \mapsto \lam{\_}{\Unit} e$, loses
    sharing. In
    $\klet x_1 = x_0 + x_0, \ldots, x_n = x_{n-1} + x_{n-1}$,
    call-by-need does linear work and call-by-name does $2^n$.
  \item \emph{Call-by-value}, $\prom{e} \mapsto e$, does work that
    call-by-need avoids: $(\slam{x} 0)\,e$ costs as much as $e$, with no
    bound relative to the source.
\end{itemize}
Call-by-need is defined by its cost: never compute what is not needed, and
never compute anything twice.

\paragraph{Time and space.}
Time needs only a count of forces and steps, which a tick monad provides
without a heap, as in Danielsson's \textsf{THUNK} monad~\cite{danielsson2008}.
Space needs the pointer structure: live space is defined by reachability, so
the semantics must know what points to what. Allocation and deallocation
behaviour is exactly what ticks cannot express.
\todo{Check Danielsson's THUNK monad: tracks time (ticks) for lazy data structures, not space.}

\paragraph{Precedents.}
The cost-sensitive view of correctness is standard:
\begin{itemize}
  \item Improvement theory for call-by-need judges program transformations
    by whether they preserve or improve time~\cite{moran1999} and
    space~\cite{gustavsson1999,gustavsson2001} up to constant factors.
    \todo{Verify Moran--Sands (time) and Gustavsson--Sands (space) and the ``up to constant factors'' formulation.}
  \item Safe-for-space compilation counts a translation correct only if it
    preserves space complexity~\cite{shao1994,paraskevopoulou2019}.
    \todo{Verify Shao--Appel citation (year, venue) and Paraskevopoulou--Appel ICFP 2019.}
  \item Okasaki's amortized bounds for lazy data structures hold only with
    memoization~\cite{okasaki1998}, and Pottier et al.\ verify such bounds
    for thunks in Iris, relying on its step-indexing~\cite{pottier2024}.
    \todo{Check the author list and venue of Thunks and Debits, and that the step-indexing claim is fair.}
\end{itemize}

\paragraph{Expressiveness and asymptotic cost.}
Closer to our question, several results separate languages by complexity
rather than by what they compute. Pippenger gives an online problem that
impure, eager Lisp solves in linear time and shows that pure, eager Lisp
needs $\Omega(n \log n)$~\cite{pippenger1997}. Bird, Jones and de Moor show
that laziness closes the gap: a lazy pure program solves the same problem
in linear time~\cite{bird1997}. That is the phenomenon we study: sharing
recovers, without observable state, efficiency that pure call-by-value
cannot reach. Hillerstr\"om, Lindley and Longley prove an analogous
separation for effect handlers: generic count takes $O(2^n)$ with
handlers and $\Omega(n2^n)$ in a pure PCF-like language~\cite{hillerstrom2020}.
Hillerstr\"om also observes that typed macro-expressibility is too coarse
to see such differences. His macro translation of deep handlers into shallow
handlers is correct and typed, yet asymptotically slower, because the
handlers it introduces accumulate~\cite{hillerstrom2021}.
\todo{Verify Pippenger (problem, bounds, venue: POPL 1996 / TOPLAS 1997) and Bird--Jones--de Moor (JFP 1997; linear-time lazy solution), cited from memory. Check the generic count bounds in Hillerstr\"om--Lindley--Longley.}

These separations are stronger than ours in one respect: they quantify over
all target programs, not only compositional translations. They are weaker
in another: they are about one problem's time, while heap adequacy
constrains every program's heap. A complexity separation for call-by-need
STLC against predicative type theory would need a cost model for type
theory, which is the problem heap adequacy avoids (\S\ref{sec:def}).

\paragraph{Reflecting cost versus preserving it.}
A semantics may \emph{reflect} cost, if its denotation computes the cost,
or a translation may \emph{preserve} cost, if running it in the host has the
same complexity. These differ:
\begin{center}
\small
\begin{tabular}{@{}lccc@{}}
  \toprule
  & reflects time & reflects space & preserves cost (CBV host) \\
  \midrule
  call-by-name thunks & no & no & no \\
  ticks~\cite{danielsson2008} & yes & no & no \\
  clairvoyant CBV~\cite{hackett2019} & yes & ? & no \\
  identity in a lazy host & no & no & only by accident \\
  semantic heap (stratified) & yes & yes & yes \\
  syntax in the heap & yes & yes & yes, with deep terms \\
  \bottomrule
\end{tabular}
\end{center}
\todo{Clairvoyant CBV (Hackett--Hutton): confirm it reflects time; does it say anything about space? Does evaluating it explore all choices?}
Only a heap gives all three. We take this as the justification for asking a
translation to keep the heap.

\paragraph{Sharing is the feature; time is the consequence.}
Both counterexamples above are failures of sharing. Call-by-name has no cell
for an argument, so it recomputes the argument at every use. A translation
that copies a cell per use has two cells where the source has one. In both
cases the extra time follows from the wrong heap. So rather than bound time
directly, we ask the translation to preserve the heap. The time savings of
call-by-need, never computing anything twice, then follow from sharing
rather than being a separate condition. This also makes the condition
checkable by a simulation proof, without a cost model for the target.

\section{Heap-adequate expressibility}
\label{sec:def}

Felleisen~\cite{felleisen1991} calls a construct $F$ of a conservative
extension $L_2 \supseteq L_1$ macro-expressible in $L_1$ when a translation
$\varphi$ is the identity on $L_1$, is homomorphic,
$\varphi(F(e_1, \ldots, e_n)) = A[\varphi(e_1), \ldots, \varphi(e_n)]$ for a
fixed $L_1$-context $A$, and preserves termination.
\todo{Check Felleisen's definitions: macro-expressibility, conservative extension, and what is preserved (termination vs observational equivalence).}
Taken directly:
\begin{itemize}
  \item Promises with an unobservable heap are macro-expressible by thunks.
    There is no gap.
  \item $\lpr$'s references are not macro-expressible in pure type theory,
    but only because state is not: $\klet x = e \kin (x, x)$ and $(e, e)$
    are equivalent in a pure language and not once $\new$ exists. This is
    equally true of ML, and has nothing to do with stratification.
\end{itemize}
Neither sees the intensional gap. Macro-expressibility is also too strict
in the wrong place. It requires the translation to fix the pure constructs,
so any feature whose translation must thread state through the whole
program fails it. Both sides of our comparison need that threading:
$\lpr$'s references do, and so do call-by-need's shared memo cells once
the heap must be preserved. Under Felleisen's definition both fail, for a
reason that has nothing to do with stratification.

We therefore leave macro-expressibility behind and define a different
notion. It keeps Felleisen's insistence that the translation be
compositional, with a fixed context per construct, but lets every construct
be reinterpreted, and it adds a heap condition.

\paragraph{Locality.} The literature distinguishes degrees of locality.
A \emph{macro-local} translation, Felleisen's original notion, is the
identity on constructs outside the feature and translates each feature
construct by a fixed context. A \emph{feature-local} translation is also
the identity outside the feature, but places no restriction on how the
feature itself is translated: a translation from named variables to de
Bruijn indices is feature-local but not macro-local. Neither admits
state-passing or CPS, which rewrite every construct~\cite{forster2017,lindley2016}.
\todo{Find and cite the source of the macro-local/feature-local terminology (quoted in notes.md); it cites Kawahara and Kameyama and Lindley for the idea. Check the Lindley citation.}
Our notion is orthogonal to both. It is weaker outside the feature, since
every construct may be reinterpreted. It is as strong as macro-locality on
the feature, since each construct must map to a fixed context, which rules
out codes (see the remark on codes below). Up to the heap condition, it asks for a
compositional, typed model of the source in the target.

\begin{definition}[Heap-adequate expressibility]
  A source language $L$ with a heap is \emph{heap-adequately
  expressible} in a type theory $\mathcal{T}$ if there is a
  translation $\varphi$ that is
  \begin{enumerate}
    \item \emph{typed}: source types map to $\mathcal{T}$ types, and
      well-typed terms to well-typed terms;
    \item \emph{compositional}: each construct maps to a fixed
      $\mathcal{T}$-context applied to the translations of its immediate
      subterms, never to their syntax. Unlike Felleisen's definition, pure
      constructs may change too, so store-passing translations are allowed;
    \item \emph{heap-adequate}: whenever a subterm is evaluated, in a
      related environment and from related heaps, the source run and the
      target computation agree in the following sense.
      \begin{itemize}
        \item They produce related results.
        \item There is an injective map from target cells to source cells,
          under which every cell's contents are related. In particular, a
          cell holding a memoized value corresponds to a cell holding a
          memoized value, not a pending computation.
        \item Every source cell that is still live, from the subterm's
          environment or its result, is the image of a target cell.
      \end{itemize}
  \end{enumerate}
\end{definition}

Heap adequacy is checked at every boundary of evaluation, not only at the
end. A closed program of type $\Nat$ ends with no live cells, so a condition
on the final heap alone would say nothing about how it got there.

\begin{remark}[Heap adequacy sees sharing]
  Call-by-name thunks fail the last condition: the source has a live cell
  for an argument, and the target has none. A translation that copies a cell
  per use fails injectivity, because two target cells map to the one source
  cell. Alternatively it keeps one cell but never memoizes it, and fails the
  contents condition. Either way, recomputation shows up in the heap before
  it shows up in time.
\end{remark}

\begin{remark}[Compositionality excludes codes]
  \label{rem:codes}
  The syntactic model is typed and heap-adequate, and its fundamental lemma
  is compositional on typing derivations. But the case for $\prom{e}$ stores
  the closure $\mathit{thunk}\;e\;\rho$, the syntax of $e$, rather than the
  denotation of $e$. It is compositional only in a target with quotation.
  The line between a shallow model and syntax in the heap is exactly the
  fixed-context condition.

  The condition is needed, not just tasteful. A feature-local translation,
  unrestricted on the feature, could send $\prom{e}$ to an interpreter
  applied to the code of $e$. Call-by-need STLC is total, so that
  interpreter exists in $\mathcal{T}$: it is the syntactic model. Without
  the condition, part~2 of the conjecture below would be false, and the
  notion would collapse to ``has a terminating interpreter''.
\end{remark}

\begin{remark}[No cost model needed]
  An earlier version of this definition bounded time and live space up to
  constant factors, under a fixed call-by-value cost model for the target.
  That needs a cost model for type theory, which has no evaluation strategy
  of its own. It also cannot be stated inside the type theory, because
  evaluation cost is not observable there. Heap adequacy avoids both
  problems. It gives up the freedom of constant factors, such as passing
  values unboxed instead of allocating cells for them. For a separation
  result, that is a fair trade.
\end{remark}
\todo{Review the switch from cost adequacy to heap adequacy, and whether the third condition is stated at the right strength.}

\section{The conjecture}
\label{sec:conj}

\begin{conjecture}
  Let $\mathcal{T}$ be a predicative type theory: Martin-L\"of type theory
  with a cumulative hierarchy of universes and strictly positive inductive
  families, whose definitions are structurally recursive.
  \begin{enumerate}
    \item $\lpr$ is heap-adequately expressible in $\mathcal{T}$.
    \item Call-by-need STLC is not.
  \end{enumerate}
\end{conjecture}
\todo{Check whether the paper already gives a shallow model in type theory.}

\paragraph{Part 1 is proved.}
Part~1 holds, with the free-variable $\lambda$ rule, and is mechanized in
Agda. Agda's type theory is an instance of $\mathcal{T}$.

\begin{theorem}[Shallow model and simulation]
  \leavevmode
  \begin{enumerate}
    \item (\texttt{LambdaPR.agda}) There is a translation of $\lpr$ typing
      derivations into Agda that is typed and compositional. Each type denotes
      an Agda type indexed by a stratified world. Each term denotes a
      store-passing Agda function over an explicit heap of semantic values,
      in which $\new$, $\deref$ and $\assign$ are cell allocation, lookup and
      update. Every definition lives in $\mathsf{Set}$ and is accepted by
      structural recursion, without postulates or fuel.
    \item (\texttt{LambdaPRAdequacy.agda}) Every semantic computation is
      simulated by $\lpr$'s operational semantics, with one flat heap and
      fresh locations as in the paper's Fig.~5. The run produces a related
      result. There is an injective map from semantic cells to runtime cells
      under which every cell's contents are related. Runtime cells outside the
      map's image are left untouched.
    \item Consequently, a closed program of type $\Nat$ evaluates
      operationally to exactly the number the model computes
      (\texttt{adequacy}), and every closed, well-typed program terminates
      (\texttt{termination}).
    \item (\texttt{boundary}) Heap adequacy holds at every boundary of
      evaluation. Run any subterm in a related environment from related
      heaps. Afterwards, the semantic heap maps injectively into the runtime
      heap with related contents. Every runtime cell still live, from the
      subterm's environment or its result, is the image of a semantic cell.
      Reachability goes directly, through closure environments, and through
      the heap.
  \end{enumerate}
\end{theorem}

The translation is typed and compositional by part~1, and heap-adequate by
parts~2 and~4, so $\lpr$ is heap-adequately expressible in Agda's type
theory. For a closed program, part~4 specializes to
\texttt{live-in-image}: the live runtime heap is no larger than the final
semantic heap, and the semantic heap is no larger than the runtime heap.

Part~4 needed one change to the value relation. The relation for functions
is behavioural, so it said nothing about what a closure captures. A related
closure must now also have every address in its environment be a cell of its
world.

\paragraph{Regions for free.}
In the model, a function at level $k$ runs its body with fresh, empty heap
layers for levels $\geq k$, and drops them when it returns. This is sound:
nothing of level $\leq k$, neither the result, nor the captured variables,
nor anything stored below level $k$, can point into those layers. So the
universe hierarchy induces region-based deallocation. The tests in
\texttt{LambdaPRTests.agda} show it: allocating a reference inside a
level-$0$ function leaves no cell in the final semantic heap, while the same
allocation at top level leaves one. The simulation relates these freed cells
to runtime cells that the operational semantics never frees. Part~4 of the
theorem shows they are garbage: no live runtime cell lies outside the map.

\subsection{Design decisions}
\label{sec:design}
\todo{Review these design decisions.}

\paragraph{Semantic heap, syntactic operational semantics.}
The model keeps semantic values in the heap. That is the point of the
exercise, and it is possible for $\lpr$ because levels stratify the heap. The
operational semantics, as in the paper, has one flat heap of syntactic values
with fresh locations. The simulation relates the two through a map from
semantic cells to runtime addresses.

\paragraph{Cells named by key.}
A semantic cell is named by its layer and its position counted from the
oldest cell in that layer. Positions counted from the newest cell shift every
time a cell is allocated, so the map would have to be updated on every
allocation. Keys never change, so the map is a plain function on keys, and
lowering a world needs no change of map.

\paragraph{The free-variable $\lambda$ rule.}
Described in \S2. With a thinning in the syntax, the semantics of $\lambda$
restricts the environment to the variables the body may use, so a
function's level reflects only what it captures.

\paragraph{Fresh layers for function bodies.}
A function at level $k$ only receives a heap of $k$ layers, but its body may
allocate at higher levels. The body therefore runs with fresh, empty upper
layers, which are dropped when it returns. This is what makes the
semantic domain definable at level $k$, and it yields the region discipline
above.

\paragraph{Behavioural relation for functions, plus what they capture.}
The relation for a function is behavioural: related arguments give related
runs. For space this is not enough, because it says nothing about the
locations a closure captures. The relation therefore also requires every
address in a closure's environment to be the address of a cell of the
function's world. This is the only change the space theorem needed.

\paragraph{Proof engineering.}
Two choices exist only to help Agda's inference:
\begin{itemize}
  \item The simulation is a function type, because it occurs inside the
    definition of the value relation. A record version would fail Agda's
    positivity check, because of the heap invariant. The proofs therefore
    use a record wrapper defined outside the mutual block.
  \item Conversions between levels are defined by recursion on the
    level proof rather than by rewriting along world equalities. That keeps
    the simulation proof free of transport.
\end{itemize}

\paragraph{Heap adequacy rather than time.}
The definition asks for heap adequacy, not a time bound (see the remarks
in \S\ref{sec:def}). Sharing is the feature that separates call-by-need from
call-by-name, and heap adequacy states it directly. A time bound states only
its consequence, and needs a cost model for the target that type theory does
not have. A time bound also cannot be stated inside the type theory, because
evaluation cost is not observable there. The internal substitute, a step
counter proved equal to the operational step count, is close to true by
construction for a compositional translation, so it would add little.

\paragraph{Remark: the shallow model's run-time cost.}
Separately from the definition, it is worth recording that the current
model would not run within a constant factor of $\lpr$ under call-by-value
evaluation. Worlds and Kripke proofs are ordinary run-time data in
\texttt{LambdaPR.agda}:
\begin{itemize}
  \item Composing world extensions does work proportional to the number of
    cells allocated.
  \item Weakening a reference recomputes its position.
  \item Most importantly, running a lower-level computation
    (\texttt{embed1}) weakens every cell in the upper layers, so a single
    call can cost time proportional to the heap.
\end{itemize}
Overheads proportional to the program's level $N$, such as casts and level
conversions, are harmless, since $N$ is fixed per program. Overheads
proportional to the heap are not. Kripke monotonicity has a run-time cost in
a shallow embedding, unless worlds are erased.

\texttt{ErasureProbe.agda} checks, in Agda~2.7 with \texttt{--erasure}, the
pieces of a representation in which each semantic operation does constant
work besides its subcalls, up to the level $N$ and list traversal for heap
lookup. List traversal costs the same as the operational semantics' heap
lookup.
\begin{itemize}
  \item Worlds and Kripke proofs are erased, so they cost nothing at run
    time.
  \item A reference is an index counted from the oldest cell of its layer,
    with an erased validity proof. Weakening it is then free, because the
    index does not change.
  \item Each heap layer records, in erased form, the world its cells were
    typed in, and a run-time count of its cells. Cells are weakened lazily,
    when read, so running a lower-level computation no longer touches the
    upper layers.
  \item Reading a cell walks the layer using run-time numbers, and fixes up
    the cell's type with an erased proof.
\end{itemize}
Porting \texttt{LambdaPR.agda} and its simulation proof to this
representation would change the value relation for references and the heap
invariant, which most of the proof depends on. Since the definition no longer
asks for a time bound, the port is optional.
\todo{Decide whether the erased representation is worth porting, as engineering, now that time is not part of the definition.}

\paragraph{Evidence for part 2: sharing forces a single $D(g)$.}
Heap adequacy forbids a translation from copying $g$'s cell or recomputing
its contents. The following family shows that this is not a technicality:
recomputation would be exponentially more expensive. Let
\[
  Q_0 = (\slam{h} h)\,(\slam{f} f\,0), \qquad
  Q_{k+1} = (\slam{r} \slam{f} r + f\,0)\,(P\,Q_k).
\]
Forcing $g = Q_{k+1}$ once allocates one promise $r$ for $P\,Q_k$, shared by
both calls of $g$'s function, so call-by-need runs $P\,Q_k$ in time linear
in $k$. A translation that copies $g$'s cell, or re-evaluates its contents,
per use runs $P\,Q_{k}$ twice at each level, in time $2^k$. A heap-adequate
translation must keep $g$'s contents as one shared semantic object, a single
target cell for the single source cell. That is all \S\ref{sec:obstruction}
needs.
\todo{Check the $Q_k$ family carefully: types, that call-by-need is linear, and that per-use copying is exponential.}

\paragraph{Toward a proof.}
The positive part is a construction. The negative part needs care, because
type theory can encode things: the syntactic model is itself a type-theoretic
definition, with codes for thunks. The proof must use compositionality to rule
out codes, not positivity alone. One route: from heap adequacy at $P$,
extract a single type $D(g)$ for the contents of $g$'s cell. That type must
appear negatively in the type of the heaps it reads. Then show that no
strictly positive, structurally recursive definition in $\mathcal{T}$ has
that shape.

\section{Discussion}

\paragraph{Codes, as for System F.}
Reynolds showed that System F has no model in which types denote
sets~\cite{reynolds1984}. Its models interpret types as relations over
untyped terms, that is, over codes. Promise heaps are similar: they have no
shallow predicative model, and the model that works uses syntax. The
difference is strength. System F's models quantify over all relations,
which predicative type theory cannot do. The syntactic model of promises
quantifies only over syntactic histories, and is a predicative definition.
\todo{Check the Reynolds 1984 claim and the PER-model description.}
\begin{center}
\small
\begin{tabular}{@{}lll@{}}
  \toprule
  & naive shallow model & what works \\
  \midrule
  System F & none~\cite{reynolds1984} & PERs over untyped terms (impredicative) \\
  promise heaps & none (conjecture) & relation over syntax (predicative) \\
  \bottomrule
\end{tabular}
\end{center}
So the impredicativity of promises is real, but small: it is tamed by
syntax. Our slogan: \emph{safe impredicative state needs intensional codes,
and $\lpr$'s levels mark the fragment that does not.}

\paragraph{Typed expressiveness gaps.}
Forster et al.\ found translations between effect constructs that exist
untyped but fail once they must preserve types, because the target lacks a
needed type~\cite{forster2017}. Ours has the same shape: an untyped
store-passing translation of promises is easy, and the obstruction is that
the heap type cannot be defined.
\todo{Check the exact claim in Forster, Kammar, Lindley and Pretnar about typed vs untyped macro-expressibility.}

\paragraph{The syntactic model.}
The model that covers all of STLC uses worlds that are histories: sequences
of stores in which each step is an allocation or a memoization, and a
memoization is justified by a run of the cell's own thunk recorded earlier
in the same world. Worlds contain no types, and the relation is defined by
induction on types. Promises cost no level. A promise is good when its
cell's contents are a good computation in every future. Acyclicity of the
heap is a consequence of the fundamental lemma, not an assumption. Details
are in \texttt{counterexample.tex}, \S8, and \texttt{Promises.agda}.

\paragraph{Consistency is not the question.}
Pure promises are consistent in type theory: Lean~4 implements memoizing
thunks whose logical meaning is a function from $\Unit$. Call-by-need STLC
is total~\cite{launchbury1993,ariola1995}. The obstruction is about the
meta-level: whether the heap can be a semantic object in a total, shallow
model. With an observable heap, Pédrot and Tabareau's fire
triangle~\cite{pedrot2020} suggests further trouble for dependent
elimination; with an unobservable heap, cost is the only reason to keep it.
\todo{Check the Lean 4 Thunk claim, and the Launchbury / Ariola et al.\ citations for termination of call-by-need STLC.}
\todo{Check the fire triangle statement: substitution, dependent elimination, observable effects.}

\paragraph{Deallocation.}
Neither model yet deallocates. Space adequacy needs reachability, so the
next step is a collection step in worlds that removes cells unreachable from
the live roots, including the continuation's environment.

\begin{thebibliography}{99}
\bibitem{ahmed2004} A.~Ahmed.
  \emph{Semantics of Types for Mutable State}. PhD thesis, Princeton, 2004.
\bibitem{ahmed2007} A.~Ahmed, M.~Fluet, G.~Morrisett.
  $L^3$: A linear language with locations. \emph{Fundamenta Informaticae},
  2007.
\bibitem{ariola1995} Z.~Ariola, M.~Felleisen, J.~Maraist, M.~Odersky,
  P.~Wadler. A call-by-need lambda calculus. POPL 1995.
\bibitem{bird1997} R.~Bird, G.~Jones, O.~de~Moor. More haste, less speed:
  lazy versus eager evaluation. \emph{J.\ Funct.\ Program.} 7(5), 1997.
\bibitem{coquand1990} T.~Coquand, C.~Paulin-Mohring. Inductively defined
  types. COLOG-88, LNCS 417, 1990.
\bibitem{danielsson2008} N.~A.~Danielsson. Lightweight semiformal time
  complexity analysis for purely functional data structures. POPL 2008.
\bibitem{felleisen1991} M.~Felleisen. On the expressive power of
  programming languages. \emph{Science of Computer Programming}, 1991.
\bibitem{forster2017} Y.~Forster, O.~Kammar, S.~Lindley, M.~Pretnar. On
  the expressive power of user-defined effects. ICFP 2017.
\bibitem{gustavsson1999} J.~Gustavsson, D.~Sands. A foundation for
  space-safe transformations of call-by-need programs. HOOTS 1999.
\bibitem{gustavsson2001} J.~Gustavsson, D.~Sands. Possibilities and
  limitations of call-by-need space improvement. ICFP 2001.
\bibitem{hackett2019} J.~Hackett, G.~Hutton. Call-by-need is clairvoyant
  call-by-value. ICFP 2019.
\bibitem{kammar2017} O.~Kammar, P.~B.~Levy, S.~K.~Moss, S.~Staton. A monad
  for full ground reference cells. LICS 2017.
\bibitem{hillerstrom2020} D.~Hillerstr\"om, S.~Lindley, J.~Longley.
  Effects for efficiency: asymptotic speedup with first-class control.
  \emph{Proc.\ ACM Program.\ Lang.} 4(ICFP):100, 2020.
\bibitem{hillerstrom2021} D.~Hillerstr\"om. \emph{Foundations for
  Programming and Implementing Effect Handlers}. PhD thesis, University of
  Edinburgh, 2021.
\bibitem{koronkevich2025} P.~Koronkevich, W.~J.~Bowman. Type universes as
  Kripke worlds. ICFP 2025.
\bibitem{launchbury1993} J.~Launchbury. A natural semantics for lazy
  evaluation. POPL 1993.
\bibitem{lindley2016} S.~Lindley. Delimited continuations, macro
  expressiveness, and effect handlers. Talk, Shonan Meeting 103, 2016.
\bibitem{moran1999} A.~Moran, D.~Sands. Improvement in a lazy context: an
  operational theory for call-by-need. POPL 1999.
\bibitem{okasaki1998} C.~Okasaki. \emph{Purely Functional Data Structures}.
  Cambridge University Press, 1998.
\bibitem{paraskevopoulou2019} Z.~Paraskevopoulou, A.~W.~Appel. Closure
  conversion is safe for space. ICFP 2019.
\bibitem{pedrot2020} P.-M.~P\'edrot, N.~Tabareau. The fire triangle: how
  to mix substitution, dependent elimination, and effects. POPL 2020.
\bibitem{pippenger1997} N.~Pippenger. Pure versus impure Lisp.
  \emph{ACM Trans.\ Program.\ Lang.\ Syst.} 19(2), 1997.
\bibitem{pottier2024} F.~Pottier, A.~Gu\'eneau, J.-H.~Jourdan, G.~M\'evel.
  Thunks and debits in separation logic with time credits. POPL 2024.
\bibitem{reynolds1984} J.~C.~Reynolds. Polymorphism is not set-theoretic.
  Semantics of Data Types, LNCS 173, 1984.
\bibitem{shao1994} Z.~Shao, A.~W.~Appel. Space-efficient closure
  representations. LFP 1994.
\end{thebibliography}
\todo{All bibliography entries are from memory: check authors, titles, venues and years.}

\end{document}
